\section{Biological Growth Models}
\index{model!biological}%

\subsection{Malthusian Model}
\index{model!Malthusian}%
\index{Malthus!Malthusian model}%
\index{exponential!Malthusian model}%

The simplest model of population growth is one in which it is assumed that the
population has sufficient nourishment and space for indefinite growth. In this
model, the function $u(x)$ represents the number of individuals at time $x$, and
it is assumed that the growth rate of the population $u'(x)$ is proportional to
the number of individuals:
\[
  u'(x) = A u(x).  
\]
The solution to this simple linear equation (exercise~\ref{ex:58230978}) is
$u(x) = u_0\cdot  e^{A x}$ where $u_0=u(0)$ is the population at the initial
time $x=0$. This means that if $A>1$, the population grows indefinitely, tending
to $+\infty$ as $x\to +\infty$ with exponential speed. If $A<1$, the population
tends to zero (i.e., becomes extinct) also with exponential speed.

\subsection{Logistic Equation}
\index{equation!logistic}%
\index{logistic!differential equation}%
\index{growth!logistic}%

To make the Malthusian model more realistic, one must consider the fact that
individuals share a common space and that the presence of many individuals in
the same space reduces resources and thus negatively contributes to population
growth. We therefore assume that the negative effect on a unit of population is
proportional to the total population: $-Bu$, while the Malthusian increase is
still given by the coefficient $A$ per unit of population. Thus, we have
\[
  u'(x)  = (A-Bu(x))\cdot u(x)
\]
This is a separable variables equation. The functions $u(x) = 0$ and $u(x) =
\frac{A}{B}$ are stationary (i.e., constant) solutions. The first represents an
empty population, which therefore can never grow. The solution $u=\frac{A}{B}$
represents a population in perfect equilibrium. To simplify the notation, we can
make the change of variable $v(x) = \frac{B}{A}\cdot u(x/A)$, which leads us to
the equivalent equation $v'=(1-v)\cdot v$. We can then assume $A=1$ and $B=1$
and return to the variable $u$:
\begin{equation}\label{eq:ode_logistica}
  u' = (1-u) u.  
\end{equation}

\newsavebox{\qrOdeLogistica}\sbox{\qrOdeLogistica}{%
\myqrshortdoclink{ode_logistica}{The solutions of the logistic equation~\eqref{eq:ode_logistica}}}
\begin{figure}
  \includegraphics[width=\textwidth]{ode_logistica.pdf}
  \caption{
The solutions of the logistic equation~\eqref{eq:ode_logistica}.
\ifwidemargin\\\\\fi%
\usebox{\qrOdeLogistica}}
\label{fig:ode_logistica}
\end{figure}

Having identified the stationary solutions $u=0$ and $u=1$, we can separate the
variables and integrate:
\[
 \int \frac{u'}{(1-u)u}\, dx = x - x_0.
\]
The integral decomposes using the method of partial fractions:
\[
  \int \frac{u'}{(1-u)u}\, dx
  = \int \frac{1}{(1-u)u}\, du 
  = \int \Enclose{\frac{1}{u} + \frac{1}{1-u}}\, du 
  = \ln \abs{u} - \ln\abs{1-u}
  = \ln \abs{\frac{u}{1-u}}.
\]
Thus, from 
\[
\abs{\frac{u}{1-u}} = e^{x-x_0}  
\]
we obtain, with some algebraic manipulation,
\[
    u(x) = 1 - \frac{1}{1\pm e^{x-x_0}}.
\]
The graph of these functions is depicted in Figure~\ref{fig:ode_logistica}. In
Chapter~\ref{sec:studio_qualitativo_edo}, we will see how it is quite simple to
obtain the qualitative behavior of the solution graphs of this equation without
having to explicitly determine the solutions.

\subsection{The Volterra-Lotka Model}
\index{model!Volterra-Lotka}%
\index{Volterra!predator-prey model}%
\index{Lotka!predator-prey model}%
\index{predator/prey}%

The Volterra-Lotka predator/prey model is the first example of a mathematical
model applied to biology. \index{Volterra Vito}%
The Italian mathematician Vito Volterra was interested in understanding, through
a mathematical model, the seemingly paradoxical phenomenon whereby the decline
in fishing activities during World War I led to a decrease in the quantity of
edible fish and a simultaneous relative increase in predatory species such as
sharks upon the resumption of fishing.

In the predator/prey model devised by Volterra, we denote by $u(x)$ the size of
the prey population and by $v(x)$ the size of the predator population as time
$x$ varies. The model is thus represented by a system of first-order nonlinear
equations:
\begin{equation}\label{eq:preda_predatore}
  \begin{cases}
    u' = A u - B u v \\
    v' = C u v - D v
  \end{cases}
\end{equation}
with $A,B,C,D$ positive real constants. The growth factor of the prey is
dictated by a Malthusian law $A u$, which represents the self-sufficiency of the
prey (it is assumed they have sufficient nourishment to grow indefinitely). The
term $-B u v$ represents the killing of prey by predators: it is a term
proportional to both the prey population $u$ and the predator population $v$ (as
if each prey had a probability proportional to $v$ of encountering a predator).
The term $C u v$ represents the corresponding growth of predators due to the
nourishment provided by the prey. Conversely, predators cannot sustain
themselves autonomously, so the term $-D v$ represents the natural extinction of
predators in the absence of nourishment.

Some solutions of the system~\eqref{eq:preda_predatore} can be found explicitly.
For example, if $v=0$, we revert to the Malthusian equation for $u$ and find
$u(x) = k e^{Ax}$: the prey grows exponentially in the absence of predators.
Similarly, if $u=0$, we find $v(x) = k e^{-Dx}$: the predators become extinct
with exponential speed in the absence of prey.

We can also find stationary (i.e., constant) solutions by setting $u'=0$ and
$v'=0$. We rewrite the system in the form:
\begin{equation}\label{eq:preda_predatore_2}
  \begin{cases}
    u' = u \cdot (A - B v) \\
    v' = v \cdot (C u - D).
  \end{cases}
\end{equation}
and observe that to simultaneously nullify $u'$ and $v'$, we must have either
$u=0$ and $v=0$ (a solution in which there are neither prey nor predators) or we
can have $u=\frac D C$ and $v=\frac A B$, which is a non-trivial equilibrium
point between the two species.

We can obtain rather precise information if we can identify an integral of the
system, that is, a function $H(u,v)$ that remains constant along the solutions:
$\frac {d}{dx} H(u(x),v(x)) = 0$. By taking the ratio of the two equations of
the system~\eqref{eq:preda_predatore_2}, we obtain
\[
  \frac{v'}{u'} 
  = \frac{v \cdot (C u - D)}{u \cdot (A - B v)}
\]
separating the variables
\[
  \frac{A - B v}{v} v' = \frac{C u - D}{u} u'
\]
integrating in $dx$
\[
  \int \frac{A - B v}{v} v'\, dx  = \int \frac{C u - D}{u} u'  \, dx
\]
making the changes of variable $u=u(x)$, $v=v(x)$, $du = u'\, dx$, $dv=v'\, dx$
and performing the divisions:
\[     
  \int \Enclose{\frac A v - B}\, dv = \int \Enclose{C  - \frac D u}\, du
\]
from which, for some constant $k$:
\[
  A \ln v - B v = C u - D\ln u + k.
\]
This means that the function 
\[
  H(u,v) = B v + C u - A \ln v - D \ln u
\]
is constant along the curves $(u(x),v(x))$. Thus, the trajectories of the system
with $u>0$ and $v>0$ are bounded, because if $u$ or $v$ were to go to infinity,
the terms $Bv-A\ln v$ or $Cu - D \ln u$ would also go to infinity, and the
function $H$ could not remain constant.

\newsavebox{\qrEdoVolterra}\sbox{\qrEdoVolterra}{%
  \myqrshortdoclink{edo_volterra}{Orbits of the predator/prey system~\ref{eq:preda_predatore}}}
\begin{figure}
  \begin{center}
    \includegraphics[width=\textwidth]{ode_volterra.pdf}
  \end{center}
  \caption{
  The orbits of the predator/prey system~\eqref{eq:preda_predatore} 
  with the choice of parameters $A=B=C=D=1$.
  \ifwidemargin\\\\\fi%
  \usebox{\qrEdoVolterra}}
  \label{fig:ode_volterra}
\end{figure}

Consider now a solution $u(x),v(x)$ that at time $x_0$ is on the segment
$u=\frac DC$ and $v\in \openinterval{0}{\frac A B}$ (refer to
Figure~\ref{fig:ode_volterra} where we have set $\frac AB=1$ and $\frac DC=1$).
Since $v<\frac A B$, we have $u'>0$, and thus the curve moves to the right. As
soon as $u>\frac DC$, we also have $v'>0$, and thus the curve moves upward.
Since $u$ cannot tend to infinity, the solution must eventually reach the line
$v=\frac AB$ where the derivative of $u$ vanishes, and the curve moves
vertically. It then becomes $v>\frac AB$ where $u'<0$. The coordinate $u$ thus
decreases, and since $v$ cannot tend to infinity, $u$ must reach the value
$\frac DC$ where $v'=0$. Then it continues with $u<\frac DC$ where $v'<0$. Thus,
both $u$ and $v$ decrease. However, they must remain positive because the
function $H$ tends to infinity also when $u\to 0$ or $v\to 0$, and the axes
$u=0$ and $v=0$ are orbits of the explicit solutions we found initially, and due
to the uniqueness of solutions, it is not possible for different orbits to meet
at a point. Thus, at some point, we will again have $v= \frac AB$ and $u\in
\openinterval{0}{\frac DC}$, and then $v<\frac AB$ with $u'>0$ and $v'<0$. Our
curve will therefore again reach at some time $x_1>x_0$ the segment from which
we started: $u=\frac DC$ and $v\in \openinterval{0}{\frac AB}$. Throughout this
path, $H$ has remained constant, so we will have
$H(u(x_1),v(x_1))=H(u(x_2),v(x_2))$. But $u(x_1)=u(x_2) = \frac DC$, and setting
$f(v) = H(D/C,v)$, we observe that
\begin{align*}
   f(v) = B v + D - A \ln v - D \ln \frac DC.
\end{align*}
studying the sign of the derivative $f'(v)=B-\frac{A}{v}$, we observe that for
$v<\frac AB$, the derivative is strictly negative, thus the function $f$ is
strictly decreasing in the interval $v\in(0,D/C)$, hence injective. Since
$f(v(x_1)) = f(v(x_2))$, we deduce that $v(x_1)=v(x_2)$. This means that the
curve closes, and the solutions are therefore periodic with period $T=x_2-x_1$.

Taking any point in the quadrant $u>0$, $v>0$ (excluding the critical point),
the solution starting from that point will follow a path that, based on the
considerations we have already made, is forced to eventually cross the segment
we considered and is, in fact, one of those we have already studied.

This qualitative study describes the behavior of the solutions of the
system~\ref{eq:preda_predatore}: try to interpret the biological significance of
what we have deduced. Numerical methods (such as those used to obtain
Figure~\ref{fig:ode_volterra}) allow us to obtain quantitative information with
the desired precision.

We can now also answer the initial question posed to Vito Volterra, namely, to
explain what happens in the presence of human fishing activity. The system
studied so far describes the behavior in the absence of fishing activity (i.e.,
what was observed by naturalists at the end of World War I). In the presence of
intensive fishing, we must assume that the quantity of fish caught is
proportional, for each species, to the size of the species itself. Thus, we must
subtract a term $Eu$ from the first equation and a corresponding term $Ev$ from
the second:
\begin{equation}
  \begin{cases}
    u' = A u - B u v - Eu = (A-E) u - B uv\\
    v' = C u v - D v - Ev = C uv - (D+E) v.
  \end{cases}
\end{equation}
We observe that the system remains of the same type, only the coefficient $A$
has decreased, becoming $A-E$, while the coefficient $D$ has increased, becoming
$D+E$. Thus, the equilibrium point $u=\frac DC$, $v=\frac AB$ has shifted to the
right and downward since $D$ has increased and $A$ has decreased. This means
that on average, fishing activity favors an increase in prey and a decrease in
predators. Conversely, the decrease in fishing activity leads to a decrease in
prey and an increase in predators.